% Einheit 3 von 4 – Grundwert berechnen (Katalog Einheit 4)
\einheitenkopf[e3][3 Grundwert]{Einheit 3 von 4 · Grundwert berechnen}
\zweigzeile{Hier lernst du, aus einem Teil und seinem Anteil in Prozent das Ganze auszurechnen · neu in diesem Jahr · P10 oft · baut auf: Prozentwert (Einheit 2), Streifen einteilen (Nr.~13), Dreisatz (Nr.~8)}
\setcounter{aufgabe}{33}

\begin{aufgabe}{Ich erkenne, was gegeben und was gesucht ist.}
Kreuze an, was gesucht ist. Rechne nicht. Das Ganze heißt Grundwert.
\begin{teile}
\teil Wie viel sind $20\,\%$ von $45\,€$?\\ \kreuz{Prozentsatz}\quad \kreuz{Prozentwert}\quad \kreuz{Grundwert}
\teil Wie viel Prozent sind 9 von 36 Kindern?\\ \kreuz{Prozentsatz}\quad \kreuz{Prozentwert}\quad \kreuz{Grundwert}
\teil $12\,€$ sind $25\,\%$ des Preises. Wie teuer ist das Spiel?\\ \kreuz{Prozentsatz}\quad \kreuz{Prozentwert}\quad \kreuz{Grundwert}
\teil Von 80 Schülern fahren $30\,\%$ mit dem Rad. Wie viele sind das?\\ \kreuz{Prozentsatz}\quad \kreuz{Prozentwert}\quad \kreuz{Grundwert}
\teil 30 Schüler fahren mit dem Rad. Das sind $40\,\%$ aller Schüler. Wie viele Schüler sind es?\\ \kreuz{Prozentsatz}\quad \kreuz{Prozentwert}\quad \kreuz{Grundwert}
\end{teile}
\end{aufgabe}

\verfahren{Das Ganze am Streifen}
\begin{aufgabe}{Ich kann am Streifen ausrechnen, wie viel das Ganze ist.}
Rechne am Streifen.

\swz{$50\,\%$ sind $10\,€$. Wie viel sind $100\,\%$?}{\streifen[20]{50}{$0\,€$}{$?$}}
\swz{$25\,\%$ sind $10\,€$. Wie viel sind $100\,\%$?}{\streifen[20]{25}{$0\,€$}{$?$}}
\swz{$20\,\%$ sind $10\,€$. Wie viel sind $100\,\%$?}{\streifen[20]{20}{$0\,€$}{$?$}}
\swz{$10\,\%$ sind $10\,€$. Wie viel sind $100\,\%$?}{\streifen[20]{10}{$0\,€$}{$?$}}
\swfrage{Was bleibt gleich, was ändert sich?}
\end{aufgabe}

\verfahren{Das Ganze über 1 \%}
\begin{aufgabe}{Ich kann über 1 \% ausrechnen, wie viel das Ganze ist.}
Rechne mit dem Dreisatz über 1 \%.
\rechenplatz{3}

\swz[1]{$6\,\%$ sind $24\,$kg. Wie viel sind $100\,\%$?}{\begin{dreisatz}{Prozent}{Kilogramm}
\dsz{6\,\%}{24\,\text{kg}}
\dsp{:6}{:6}
\dsz{1\,\%}{\dsleer}
\dsp{\cdot 100}{\cdot 100}
\dsz{100\,\%}{\dsleer}
\end{dreisatz}}
\swz[1]{$3\,\%$ sind $1{,}20\,€$. Wie viel sind $100\,\%$?}{\begin{dreisatz}{Prozent}{Euro}
\dsz{3\,\%}{1{,}20\,€}
\dsp{\phantom{:0}}{\phantom{:0}}
\dsz{1\,\%}{\dsleer}
\dsp{\phantom{:0}}{\phantom{:0}}
\dsz{100\,\%}{\dsleer}
\end{dreisatz}}
\anweisung{Rechne mit dem Taschenrechner.}
\swz[1]{$37\,\%$ sind $81{,}40\,€$. Wie viel sind $100\,\%$?}{\begin{dreisatz}{Prozent}{Euro}
\dsz{37\,\%}{81{,}40\,€}
\dsp{\phantom{:0}}{\phantom{:0}}
\dsz{1\,\%}{\dsleer}
\dsp{\phantom{:0}}{\phantom{:0}}
\dsz{100\,\%}{\dsleer}
\end{dreisatz}}
\end{aufgabe}

\verfahren{Das Ganze im Sachtext}
\begin{aufgabe}{Ich kann in einem Sachtext ausrechnen, wie viel das Ganze ist.}
Berechne das Ganze.

\swz[3]{In der Klasse 7b fahren 12 Kinder mit dem Bus. Das sind $40\,\%$ der Klasse. Wie viele Kinder sind in der Klasse?}{\streifen{40}{$0$}{$?$}}
\swz[3]{In einem Bus sind $25\,\%$ der Plätze besetzt. 45 Plätze sind frei. Wie viele Plätze hat der Bus?}{\streifen[4]{25}{$0$}{$?$}}
\end{aufgabe}

\begin{aufgabe}{Ich kann unterscheiden, ob der Teil, der Anteil in Prozent oder das Ganze gesucht ist.}
Entscheide zuerst, was gesucht ist. Rechne dann.

\swz[3]{Wie viel sind $30\,\%$ von $50\,$kg?}{\streifen{0}{$0\,\%$}{$100\,\%$}}
\swz[3]{$18\,€$ sind $30\,\%$ des Preises. Wie hoch ist der Preis?}{\streifen{0}{$0\,\%$}{$100\,\%$}}
\swz[3]{9 von 45 Kindern tragen eine Brille. Wie viel Prozent sind das?}{\streifen{0}{$0\,\%$}{$100\,\%$}}
\swz[3]{Beim Kauf einer Mütze spart Ole $4\,€$. Das sind $25\,\%$ des alten Preises. Wie viel kostete die Mütze vorher? (P10 2025 OS)}{}
\end{aufgabe}

\begin{aufgabe}{Ich finde den Fehler, wenn jemand das Ganze ausrechnet.}
Finde den Fehler, benenne ihn und rechne richtig.

\swz[3]{$15\,€$ sind $25\,\%$ des Preises. Wie hoch ist der Preis? Tim rechnet so:}{\rechnung{0{,}25 \cdot 15\,€ &= 3{,}75\,€}}
\swz[3]{36 Kinder sagen ja. Das sind $60\,\%$ aller Befragten. Wie viele Kinder wurden befragt? Paul rechnet so:}{\rechnung{0{,}6 \cdot 36 &= 21{,}6 \\ \text{Befragte} &\approx 22}}
\end{aufgabe}

\begin{aufgabe}{Ich kann begründen, warum man für das Ganze mal 100 rechnet.}
Begründe.

\swfrage{Warum rechnet man den Wert für $1\,\%$ mal 100, um das Ganze zu bekommen?}
\swfrage{Ist das Ganze größer oder kleiner als der Teil, wenn der Teil $25\,\%$ ist? Und wenn der Teil $125\,\%$ ist?}
\end{aufgabe}

\clearpage
\begin{aufgabe}{Ich kann mit dem Ganzen planen, wie viele Plätze gebraucht werden.}
Für die Klassenfahrt haben sich schon 54 Kinder angemeldet. Das sind $60\,\%$ aller Kinder des Jahrgangs. Ein Bus hat 50 Plätze. Rechne und entscheide.

\swz[3]{Wie viele Kinder hat der Jahrgang?}{\streifen{60}{$0$}{$?$}}
\swfrage{Reichen zwei Busse, wenn alle Kinder des Jahrgangs mitfahren? Begründe.}
\end{aufgabe}

\uebersichtskasten{Grundwert: erst 1 \% (Prozentwert : p), dann mal 100. \\ 36 € sind 45 \%: \quad 1 \% $= 0{,}80$ € $\rightarrow$ 100 \% $= 80$ € \qquad Streifen: 45 \% $\widehat{=}$ 36 €, 100 \% $\widehat{=}$ ? \\ Erst zuordnen: Was ist gegeben (W, p, G), was gesucht? Dann rechnen.}
